\chapter{Lezione 20 --- Controllo ottimo LQ continuo, equazioni di Riccati e valori singolari}

\section{Dal cost-to-go all'equazione di Bellman nel continuo}

Consideriamo un sistema dinamico non lineare in tempo continuo
\[
\dot x = f(x,u),
\]
e un indice di ottimalità su orizzonte infinito
\[
J=\int_{t_0}^{+\infty}\ell(x(\tau),u(\tau))\,d\tau.
\]

Definiamo il \emph{cost-to-go} ottimo
\[
\psi(x_0)=\min_{u(\cdot)} J
\qquad\text{con } x(t_0)=x_0.
\]

Ragionando come nella programmazione dinamica, il costo ottimo a partire da \(x(t)\) deve soddisfare
\[
\psi(x(t))
=
\min_{u(\cdot)}
\left[
\int_t^{t+h}\ell(x(\tau),u(\tau))\,d\tau
+
\psi(x(t+h))
\right].
\]

Portando \(\psi(x(t))\) a destra, dividendo per \(h\) e facendo tendere \(h\to 0^+\), si ottiene
\[
\min_u\left[\dot\psi(x(t))+\ell(x(t),u(t))\right]=0.
\]

Usando la regola della catena,
\[
\dot\psi(x(t))
=
\nabla_x \psi(x(t))^T \dot x(t)
=
\nabla_x \psi(x(t))^T f(x(t),u(t)),
\]
segue l'\textbf{equazione di Bellman}:
\[
\boxed{
\min_u\left[\nabla_x\psi(x)^T f(x,u)+\ell(x,u)\right]=0.
}
\]

\section{Caso LQ continuo su orizzonte infinito}

Consideriamo ora il caso lineare-quadratico:
\[
\dot x = Ax+Bu,
\]
con funzionale di costo
\[
J=\frac12\int_0^{+\infty}\left(x^TQx+u^TRu\right)\,d\tau,
\]
dove
\[
Q=Q^T\succeq 0,
\qquad
R=R^T\succ 0.
\]

Cerchiamo una soluzione del cost-to-go nella forma quadratica
\[
\psi(x)=\frac12 x^T P x,
\qquad
P=P^T.
\]

Poiché \(P\) è simmetrica,
\[
\nabla_x\psi(x)=Px.
\]

Sostituendo nell'equazione di Bellman si ottiene
\[
\min_u
\left[
x^TP(Ax+Bu)+\frac12 x^TQx+\frac12 u^TRu
\right]=0.
\]

\subsection{Calcolo del controllo ottimo}

Per trovare il minimo rispetto a \(u\), imponiamo la condizione del primo ordine:
\[
\nabla_u\!\left(
x^TPBu+\frac12 u^TRu
\right)=0.
\]

Essendo
\[
\nabla_u(x^TPBu)=B^TPx,
\qquad
\nabla_u\!\left(\frac12 u^TRu\right)=Ru,
\]
segue
\[
B^TPx+Ru=0,
\]
cioè
\[
\boxed{
u^\star=-R^{-1}B^TPx.
}
\]

Quindi il controllo ottimo è una retroazione lineare dello stato, con guadagno
\[
K=R^{-1}B^TP.
\]

\section{Equazione algebrica di Riccati nel continuo}

Sostituiamo \(u^\star\) nell'equazione di Bellman:
\[
x^TPAx
-\,
x^TPBR^{-1}B^TPx
+
\frac12 x^TQx
+
\frac12 x^TPBR^{-1}B^TPx
=0.
\]

Riordinando:
\[
x^TPAx
+\frac12 x^TQx
-\frac12 x^TPBR^{-1}B^TPx
=0.
\]

Ora osserviamo che lo scalare \(x^TPAx\) può essere simmetrizzato:
\[
x^TPAx
=
\frac12 x^T(PA+A^TP)x.
\]

Otteniamo dunque
\[
\frac12 x^T
\left(
A^TP+PA-PBR^{-1}B^TP+Q
\right)x=0
\qquad \forall x.
\]

Perciò la matrice tra parentesi deve annullarsi:
\[
\boxed{
A^TP+PA-PBR^{-1}B^TP+Q=0.
}
\]

Questa è l'\textbf{equazione algebrica di Riccati a tempo continuo}.

\subsection{Sistema in ciclo chiuso}

Ponendo
\[
u^\star=-Kx,
\qquad
K=R^{-1}B^TP,
\]
la dinamica in ciclo chiuso diventa
\[
\dot x=(A-BK)x.
\]

Quindi la matrice di stato del sistema controllato ottimamente è
\[
A_{cc}=A-BR^{-1}B^TP.
\]

\section{Stabilità e legame con Lyapunov}

Partiamo dall'equazione di Riccati
\[
A^TP+PA-PBR^{-1}B^TP+Q=0.
\]

Poiché
\[
K=R^{-1}B^TP,
\]
si ha
\[
A_{cc}=A-BK=A-BR^{-1}B^TP.
\]

Sostituendo, l'equazione di Riccati può essere riscritta come
\[
(A-BK)^TP + P(A-BK) + \widetilde Q = 0,
\]
dove
\[
\widetilde Q = Q + PBR^{-1}B^TP \succeq 0.
\]

Cioè
\[
\boxed{
A_{cc}^TP + PA_{cc} = -\widetilde Q.
}
\]

Questa è una \textbf{equazione di Lyapunov} per il sistema in ciclo chiuso.

Se \(P\succ 0\) e \(\widetilde Q \succ 0\), allora la funzione
\[
V(x)=\frac12 x^TPx
\]
è una funzione di Lyapunov e il sistema in ciclo chiuso è asintoticamente stabile.

\paragraph{Interpretazione}
La soluzione \(P\) dell'equazione di Riccati non serve solo a scrivere il controllo ottimo:
essa codifica anche una funzione energia/costo che certifica la stabilità del sistema controllato.

\section{Ipotesi standard: stabilizzabilità e rilevabilità}

Negli appunti compare la seguente idea fondamentale.

\begin{itemize}
    \item \textbf{Stabilizzabile}: gli autovalori non controllabili di \(A\) hanno parte reale strettamente negativa.
    \item \textbf{Rilevabile} (detectable): gli autovalori non osservabili hanno parte reale strettamente negativa.
\end{itemize}

Per il problema LQ continuo, sotto le ipotesi standard
\[
(A,B)\ \text{stabilizzabile},
\qquad
(A,\sqrt Q)\ \text{rilevabile},
\]
l'equazione algebrica di Riccati ammette una \emph{soluzione stabilizzante} \(P=P^T\succeq 0\), e il feedback
\[
u^\star=-R^{-1}B^TPx
\]
rende stabile il sistema in ciclo chiuso.

\medskip

\noindent
\textbf{Nota.}
Negli appunti la conclusione è espressa in forma forte come esistenza di una soluzione \(P>0\).
Per essere più precisi, la formulazione standard parla di soluzione simmetrica positiva semidefinita stabilizzante; nei casi ``ben posti'' trattati a lezione essa è spesso positiva definita.

\section{Caso LQ continuo su orizzonte finito}

Se invece considero un orizzonte finito \([0,t_f]\), il funzionale di costo diventa
\[
J=
\int_0^{t_f}
\left(
\frac12 x^TQx+\frac12 u^TRu
\right)d\tau
+
\frac12 x^T(t_f)Sx(t_f),
\]
con
\[
S=S^T\succ 0.
\]

In questo caso il cost-to-go dipende anche dal tempo:
\[
\psi(x,t)=\frac12 x^T P(t)x.
\]

La legge di controllo ottimo conserva la stessa forma
\[
\boxed{
u^\star(t)=-R^{-1}B^TP(t)x(t),
\qquad 0<t<t_f,
}
\]
ma il guadagno non è più costante: dipende dal tempo tramite \(P(t)\).

La matrice \(P(t)\) soddisfa l'\textbf{equazione differenziale di Riccati}:
\[
\boxed{
-\dot P(t)=A^TP(t)+P(t)A-P(t)BR^{-1}B^TP(t)+Q,
}
\]
con condizione terminale
\[
\boxed{
P(t_f)=S.
}
\]

\paragraph{Osservazione sul segno}
L'equazione è scritta con il segno ``\(-\dot P\)'' perché la si integra all'indietro, partendo dal tempo finale \(t_f\) e tornando verso \(0\).  
Negli appunti questa dipendenza all'indietro è il punto concettuale importante.

\section{Richiamo di algebra lineare: decomposizione ai valori singolari}

Nella seconda parte della lezione compare un richiamo sui \textbf{valori singolari}, utile per ragionare sul condizionamento numerico delle matrici.

Sia
\[
M\in\mathbb R^{m\times k},
\qquad m\ge k.
\]
Allora esistono matrici ortogonali
\[
U\in\mathbb R^{m\times m},
\qquad
V\in\mathbb R^{k\times k},
\]
tali che
\[
U^TU=I,
\qquad
V^TV=I,
\]
e una matrice \(\Sigma\in\mathbb R^{m\times k}\) della forma
\[
\Sigma=
\begin{bmatrix}
\Sigma'\\
0
\end{bmatrix},
\qquad
\Sigma'=
\operatorname{diag}(\sigma_1,\sigma_2,\dots,\sigma_k),
\]
con
\[
\sigma_1\ge \sigma_2\ge \cdots \ge \sigma_k\ge 0,
\]
tale che
\[
\boxed{
M=U\Sigma V^T.
}
\]

Questa è la \textbf{SVD} (Singular Value Decomposition).

\subsection{Interpretazione dei valori singolari}

I numeri \(\sigma_i\) sono i \textbf{valori singolari} di \(M\).
Possono essere visti come una generalizzazione degli autovalori al caso di matrici non necessariamente quadrate.

Essi misurano anche quanto la matrice sia vicina alla singolarità.

\begin{itemize}
    \item Se \(M\) è quadrata e invertibile, allora il più piccolo valore singolare è positivo:
    \[
    \sigma_k>0.
    \]
    \item Se \(\sigma_k=0\), la matrice perde rango.
\end{itemize}

\subsection{Numero di condizionamento}

La distanza dalla non invertibilità si misura tramite il rapporto
\[
\kappa(M)=\frac{\sigma_1}{\sigma_k},
\]
detto \textbf{numero di condizionamento}.

\begin{itemize}
    \item Se
    \[
    \frac{\sigma_1}{\sigma_k}\approx 1,
    \]
    la matrice è ben condizionata e lontana dalla singolarità.
    \item Se
    \[
    \frac{\sigma_1}{\sigma_k}\gg 1,
    \]
    la matrice è male condizionata, cioè quasi non invertibile.
\end{itemize}

\paragraph{Perché compare qui?}
Nel controllo ottimo e, più in generale, nei problemi di calcolo numerico del controllo, si invertono spesso matrici come
\[
R,\qquad P+B^TT_{k+1}B,\qquad \text{ecc.}
\]
Perciò il condizionamento numerico è importante: una matrice formalmente invertibile può essere problematica dal punto di vista computazionale se è molto vicina alla singolarità.

\section{Sintesi finale della lezione}

In questa lezione abbiamo visto che:

\begin{enumerate}
    \item dal cost-to-go nel continuo si ottiene l'equazione di Bellman
    \[
    \min_u\left[\nabla_x\psi(x)^Tf(x,u)+\ell(x,u)\right]=0;
    \]
    \item nel caso LQ continuo, cercando \(\psi(x)=\tfrac12 x^TPx\), si ricava il controllo ottimo
    \[
    u^\star=-R^{-1}B^TPx;
    \]
    \item la matrice \(P\) deve soddisfare l'equazione algebrica di Riccati
    \[
    A^TP+PA-PBR^{-1}B^TP+Q=0;
    \]
    \item tale equazione si collega a una Lyapunov del sistema in ciclo chiuso
    \[
    A_{cc}^TP+PA_{cc}=-\widetilde Q;
    \]
    \item su orizzonte finito compare invece l'equazione differenziale di Riccati
    \[
    -\dot P=A^TP+PA-PBR^{-1}B^TP+Q,
    \qquad
    P(t_f)=S;
    \]
    \item i valori singolari permettono di capire se una matrice è ben o mal condizionata, tramite
    \[
    \kappa(M)=\frac{\sigma_1}{\sigma_k}.
    \]
\end{enumerate}